% GENERATED FILE. Edit canonical lessons, then run npm run generate:books.
% canonical-source-hash: fnv1a64:296c6b0dee7b6b35
% canonical-lessons: IT-C318-invito, IT-C318-invitato, IT-C318-ricevimento, IT-C318-rinfresco, IT-C318-confetto

\chapter{Gli invitati}
\label{ch:it-gli-invitati}
\hlchaptermodality{318}

\begin{chapteropening}
\textbf{By the end of this chapter:} \emph{I can talk about guests and invitations, the reception, the refreshments and the sugared almonds.}

Complete the last lesson of chapter 318: I can talk about guests and invitations, the reception, the refreshments and the sugared almonds.
\end{chapteropening}

\section[l'invito]{l'invito — the invitation}

\label{lesson:IT-C318-invito}



\begin{quote}
Before the new one: say the Italian for the corkscrew, then the Italian for the banquet.
\end{quote}

\subsection*{You'll want to know: l'invito}
\textbf{l'invito} — \textquotedblleft{}the invitation\textquotedblright{}.

\begin{quote}
Grazie per l'invito!
\end{quote}

\begin{quote}
\emph{Thank you for the invitation!}
\end{quote}

Built on \emph{invitare}.

\begin{cousinweb}
\emph{invito} is built on \emph{invitare} \textquotedblleft{}to invite\textquotedblright{}, from Latin \emph{invītāre} (Lewis \& Short).
\end{cousinweb}

\begin{grammarlens}[title={What you've built}]
One more word for telling, explaining and describing at B1.
\end{grammarlens}

\subsection*{Your turn}
\begin{itemize}
\raggedright
  \item \emph{Say it:} \emph{l'invito}
  \item \emph{Say it:} \emph{l'invito}, once more
  \item \emph{Say it:} a sentence of your own with \emph{l'invito}
  \item \emph{Recall:} say the Italian for the corkscrew, then the Italian for the banquet, then say \emph{l'invito} again
\end{itemize}

\begin{tcolorbox}[breakable,colback=teal!4,colframe=teal!35!black,title={Before you move on}]
What does l'invito mean? (\textquotedblleft{}The invitation\textquotedblright{}.) Say it once more.
\end{tcolorbox}
\section[l'invitato]{l'invitato — the guest}

\label{lesson:IT-C318-invitato}



\begin{quote}
Before the new one: say the Italian for the banquet, then the Italian for the invitation.
\end{quote}

\subsection*{You'll want to know: l'invitato}
\textbf{l'invitato} — \textquotedblleft{}the guest\textquotedblright{}.

\begin{quote}
Al matrimonio ci sono duecento invitati.
\end{quote}

\begin{quote}
\emph{There are two hundred guests at the wedding.}
\end{quote}

The past participle of \emph{invitare} used as a noun.

\begin{cousinweb}
\emph{invitato} is the past participle of \emph{invitare}, from Latin \emph{invītāre} \textquotedblleft{}to invite\textquotedblright{} (Lewis \& Short).
\end{cousinweb}

\begin{grammarlens}[title={What you've built}]
One more word for telling, explaining and describing at B1.
\end{grammarlens}

\subsection*{Your turn}
\begin{itemize}
\raggedright
  \item \emph{Say it:} \emph{l'invitato}
  \item \emph{Say it:} \emph{l'invitato}, once more
  \item \emph{Say it:} a sentence of your own with \emph{l'invitato}
  \item \emph{Recall:} say the Italian for the banquet, then the Italian for the invitation, then say \emph{l'invitato} again
\end{itemize}

\begin{tcolorbox}[breakable,colback=teal!4,colframe=teal!35!black,title={Before you move on}]
What does l'invitato mean? (\textquotedblleft{}The guest\textquotedblright{}.) Say it once more.
\end{tcolorbox}
\section[il ricevimento]{il ricevimento — the reception}

\label{lesson:IT-C318-ricevimento}



\begin{quote}
Before the new one: say the Italian for the invitation, then the Italian for the guest.
\end{quote}

\subsection*{You'll want to know: il ricevimento}
\textbf{il ricevimento} — \textquotedblleft{}the reception\textquotedblright{}.

\begin{quote}
Il ricevimento è in un giardino in collina.
\end{quote}

\begin{quote}
\emph{The reception is in a garden in the hills.}
\end{quote}

In a hotel, the front desk is also \emph{il ricevimento}.

\begin{cousinweb}
\emph{ricevimento} is built on \emph{ricevere} \textquotedblleft{}to receive\textquotedblright{}, from Latin \emph{recipere} (Lewis \& Short).
\end{cousinweb}

\begin{grammarlens}[title={What you've built}]
One more word for telling, explaining and describing at B1.
\end{grammarlens}

\subsection*{Your turn}
\begin{itemize}
\raggedright
  \item \emph{Say it:} \emph{il ricevimento}
  \item \emph{Say it:} \emph{il ricevimento}, once more
  \item \emph{Say it:} a sentence of your own with \emph{il ricevimento}
  \item \emph{Recall:} say the Italian for the invitation, then the Italian for the guest, then say \emph{il ricevimento} again
\end{itemize}

\begin{tcolorbox}[breakable,colback=teal!4,colframe=teal!35!black,title={Before you move on}]
What does il ricevimento mean? (\textquotedblleft{}The reception\textquotedblright{}.) Say it once more.
\end{tcolorbox}
\section[il rinfresco]{il rinfresco — the refreshments, the drinks party}

\label{lesson:IT-C318-rinfresco}



\begin{quote}
Before the new one: say the Italian for the guest, then the Italian for the reception.
\end{quote}

\subsection*{You'll want to know: il rinfresco}
\textbf{il rinfresco} — \textquotedblleft{}the refreshments, the drinks party\textquotedblright{}.

\begin{quote}
Dopo il concerto c'è un piccolo rinfresco.
\end{quote}

\begin{quote}
\emph{After the concert there are some refreshments.}
\end{quote}

Built on \emph{fresco}, ‘fresh’.

\begin{cousinweb}
\emph{rinfresco} is built on \emph{rinfrescare} \textquotedblleft{}to refresh\textquotedblright{}, from \emph{fresco} \textquotedblleft{}fresh\textquotedblright{}, which EtymDB traces through Medieval Latin \emph{friscus} to a Germanic word.
\end{cousinweb}

\begin{grammarlens}[title={What you've built}]
One more word for telling, explaining and describing at B1.
\end{grammarlens}

\subsection*{Your turn}
\begin{itemize}
\raggedright
  \item \emph{Say it:} \emph{il rinfresco}
  \item \emph{Say it:} \emph{il rinfresco}, once more
  \item \emph{Say it:} a sentence of your own with \emph{il rinfresco}
  \item \emph{Recall:} say the Italian for the guest, then the Italian for the reception, then say \emph{il rinfresco} again
\end{itemize}

\begin{tcolorbox}[breakable,colback=teal!4,colframe=teal!35!black,title={Before you move on}]
What does il rinfresco mean? (\textquotedblleft{}The refreshments, the drinks party\textquotedblright{}.) Say it once more.
\end{tcolorbox}
\section[il confetto]{il confetto — the sugared almond}

\label{lesson:IT-C318-confetto}



\begin{quote}
Before the new one: say the Italian for the reception, then the Italian for the refreshments.
\end{quote}

\subsection*{You'll want to know: il confetto}
\textbf{il confetto} — \textquotedblleft{}the sugared almond\textquotedblright{}.

\begin{quote}
A ogni invitato diamo cinque confetti.
\end{quote}

\begin{quote}
\emph{We give each guest five sugared almonds.}
\end{quote}

White for weddings, pink or blue for a birth. In Italian \emph{confetti} are sweets, not paper.

\begin{cousinweb}
\emph{confetto} is inherited from Latin \emph{cōnfectus} \textquotedblleft{}made, prepared\textquotedblright{} (EtymDB), from \emph{cōnficere}.
\end{cousinweb}

\begin{grammarlens}[title={What you've built}]
One more word for telling, explaining and describing at B1.
\end{grammarlens}

\subsection*{Your turn}
\begin{itemize}
\raggedright
  \item \emph{Say it:} \emph{il confetto}
  \item \emph{Say it:} \emph{il confetto}, once more
  \item \emph{Say it:} a sentence of your own with \emph{il confetto}
  \item \emph{Recall:} say the Italian for the reception, then the Italian for the refreshments, then say \emph{il confetto} again
\end{itemize}

\begin{tcolorbox}[breakable,colback=teal!4,colframe=teal!35!black,title={Before you move on}]
What does il confetto mean? (\textquotedblleft{}The sugared almond\textquotedblright{}.) Say it once more.
\end{tcolorbox}
